We approach the implementation of LLVManim as a 3-layer architecture:
\textbf{Ingestion}, \textbf{Transformation}, and \textbf{Presentation}. The Ingestion layer collects LLVM IR and metadata. The transformation layer maps the collected IR constructs and trace events to an intermediate scene graph. The presentation layer renders the scene graph with Manim animations and exposes a small API for interactive playback.

\begin{figure}[htbp]
\centerline{\includegraphics[width=\linewidth]{ModelDiagram}}
\caption{Model diagram illustrating the three-layer architecture of LLVManim. The Ingestion layer collects LLVM IR and metadata, the Transformation layer maps it to an intermediate scene graph, and the Presentation layer renders it with Manim animations.}
\label{fig:model_diagram}
\end{figure}

\subsubsection{Presentation Layer}
The Presentation layer renders the scene graph using Manim and provides a simple interface for the user to render \& view animations.


\begin{itemize}
    \item \textbf{User Interface:} Command-line interface with automatic compilation:
    \begin{lstlisting}[language=bash,caption={Mock CLI commands},gobble=4]
    # Visualize C program
    $ llvmanim demo.c --mode=cfg --output=demo.mp4
    # Custom speed, colors, & output format
    $ llvmanim demo.c --mode=memory --speed=0.75 --palette=dark --output=demo.gif
    # Interactive mode
    $ llvmanim demo.c --mode=trace --interactive
    \end{lstlisting}
    The CLI is intended for both developers integrating LLVManim into existing workflows and educators preparing course material. When the tool encounters unsupported IR constructs, missing debug metadata, or malformed input, it will emit a descriptive error message and continue processing the remainder of the program where possible.
    \item \textbf{Customization API:} Provide configuration options using command line flags and YAML files:
    \begin{itemize}
        \item \texttt{--speed <factor>}: Animation playback speed
        \item \texttt{--palette <name>}: Color scheme (dark, light, colorblind)
        \item \texttt{--mode <type>}: Visualization mode (cfg, memory, trace)
        \item \texttt{--verbose}: Include IR instruction text
    \end{itemize}
    \item \textbf{Manim Code Generation:} Translate scene graph into Python code defining a Manim \texttt{Scene} class. Visual elements become \texttt{Mobject} instances with animation commands.
    \item \textbf{Rendering Pipeline:} Invoke Manim to produce supported output formats (MP4, MOV, WEBM, \& GIF) or enable interactive openGL preview with keyboard controls.
\end{itemize}

\subsubsection{Transformation Layer}
The Transformation layer maps LLVM constructs to an intermediate scene graph for Manim rendering.

\begin{itemize}
    \item \textbf{Instruction-to-Animation Mapping:} Define animation primitives for each LLVM instruction:
    \begin{itemize}
        \item \texttt{alloca}: Create stack frame with variable label
        \item \texttt{load}/\texttt{store}: Animate arrow to/from memory location
        \item \texttt{call}: Push new stack frame with fade-in
        \item \texttt{ret}: Pop stack frame with fade-out
        \item \texttt{br}: Highlight control flow edge with color change
    \end{itemize}

    \item \textbf{Scene Graph Construction:} Create hierarchical representation where nodes correspond to visual elements (basic blocks, stack frames, memory cells) with properties for position, appearance, and timing.
    
    \item \textbf{Input:} Structured representation of program control flow and memory operations from the Ingestion layer.
    
    \item \textbf{Output (Application Facade):} Generated Manim instructions, sent to the Presentation layer for rendering.

    \item \textbf{(Optional) Timing Coordination:} Use parsed profiling data to calculate timing offsets for animation events (user-configurable).
\end{itemize}


\subsubsection{Ingestion Layer}
The Ingestion layer extracts program structure from LLVM IR. We use the \texttt{LLVMlite} \cite{llvmlite} Python binding to parse .ll files generated by clang with the \texttt{-g} flag for debug metadata.

\begin{itemize}
    \item \textbf{Static Data Collection:} \begin{itemize}
      \item \textit{Static Data Extraction:} Load LLVM modules and extract functions, basic blocks, instructions, and metadata using \texttt{LLVMlite}.
      \item \textit{Control Flow Extraction:} Use LLVM built-in analysis passes (e.g. \texttt{-opt}) to extract control-flow data from the IR.
    \end{itemize}

    \item \textbf{Input:} Source files (automatically compiled to LLVM IR) or .ll files directly.

    \item \textbf{Output:} Structured representation of the program's control flow and memory operations for the Transformation layer to convert into a Manim scene.
    
    \item \textbf{(Optional) Dynamic Data Collection:} Extract instruction timing data from LLVM profiling tools, such as \texttt{LLVM-XRay} \cite{llvm_xray}, to enable performance-based animation timing.
\end{itemize}

LLVManim will target a well-defined subset of LLVM IR, including standard control flow constructs (branches and loops), common memory operations (\texttt{alloca}, \texttt{load}, and \texttt{store}), and basic arithmetic and comparison instructions. Unsupported constructs like intrinsics, exception handling, or vector operations will be skipped with a warning message indicating which instructions were omitted, keeping the tool usable even on programs that use unsupported features.

\begin{figure}[htbp]
\centerline{\includegraphics[width=\linewidth]{MockOut}}
\caption{Mock output of LLVManim showing a basic increment operation.}
\label{fig:mock_output}
\end{figure}
